\scanpage{21}
\uline{Re NB.} Soit tjrs $g=\dim\mathrm{rel.}$ de $A/S$, et considérons
\[
\omega_A=\deleted{\unsure{A}}\operatorname{dét}(\underline t_A)^{-1}
=\bigwedge^g\underline{\unsure{t|\omega}}_A\deleted{\unsure{\vee}}
\]
\inserted{(forme diff. de degré max le long de la section unité)}
C'est un faisceau inv. à $p$-structure non dégénérée, i.e. déduit par ext. des scalaires d'un \deleted{\unsure{faisceau}}\inserted{\deleted{\unsure{fais}} faisceau loc. ct} \deleted{$\unsure{\mathbb F_p\text{-vectoriel}}$} de rang un sur $\mathbb F_p$, savoir $\deleted{\unsure{(}}\operatorname{dét}({}_pA^*)_{\text{et}}\deleted{\unsure{)}}$. À ce titre, $\omega_A^{\otimes(p-1)}$ admet une section canonique.

Si on ne suppose plus var. $A/S$ ordinaire, et si $U$ est l'ouvert des $s\in S$ tels que $A_s$ soit ordinaire, \deleted{\unsure{tt}} ce qui précède s'applique à $A/U$, en particulier on trouve une section canonique \inserted{$\sigma$} de $L\vert U$, partout $\ne0$.
\deleted{Si $U$ est sch. dense dans $S$, on peut considérer $\sigma$ comme une \deleted{\unsure{par}} pseudo\deleted{\unsure{fonction}}\inserted{pluri-forme} sur $S$ (ou section méromorphe sur $S$), \deleted{\unsure{qui se}}\inserted{mais en}}

Mais en fait, $\sigma$ provient d'une section canonique de $L^{\otimes(p-1)}$ sur $S$ tt entier, définie à l'aide de
\diagram{\begin{tikzcd}[column sep=large]t_A^{(p)}\arrow[r,"\pi"] & t_A\end{tikzcd}}
en passant au dét, d'où
\diagram{\begin{tikzcd}[column sep=large](\operatorname{dét}\underline t_A)^{\otimes p}\arrow[r] & \operatorname{dét}\underline t_A\end{tikzcd}}
d'où
\[
\sigma_A\in\Gamma(\operatorname{dét}\underline t_A)^{1-p}=\Gamma\omega_A^{(p-1)}.
\]
\mspage{20}
